\documentclass[11pt, oneside]{article} 
\usepackage{geometry}
\geometry{letterpaper} 
\usepackage{graphicx}
	
\usepackage{amssymb}
\usepackage{amsmath}
\usepackage{parskip}
\usepackage{color}
\usepackage{hyperref}

\graphicspath{{/Users/telliott/Dropbox/Github-Math/figures/}}
% \begin{center} \includegraphics [scale=0.4] {gauss3.png} \end{center}

\title{Pythagorean Theorem:  Pappus}
\date{}

\begin{document}
\maketitle
\Large

%[my-super-duper-separator]

Pappus came up with a beautiful result which includes the Pythagorean theorem as a special case.  First, we need to recall a result about parallelograms, Euclid I.35.

\begin{center} \includegraphics [scale=0.15] {EI_35.png} \end{center}

If two parallelograms lie on the same base and with the opposite side on the same parallel, then they are equal in area.

Here, the base $BC$ is shared.  However, sequential applications of the same theorem will show that the result obtains, even if neither of the opposing sides is shared.

\subsection*{Pappus's parallelogram theorem}

\label{sec:PProof_Pappus}

In the figure below, on the two sides of any $\triangle ABC$ draw any parallelograms $ACDE$ and $BCFG$.  Extend the two new sides to meet at $H$.

\begin{center} \includegraphics [scale=0.22] {Pappus_pgram1.png} \end{center}

Draw $HC$ and its extension such that it cuts $AB$ at $T$ and then let $HC = TU$.

Draw $AJ \parallel HCTU \parallel BI$ and $JUI \parallel ATB$.

\emph{Proof}.

First, $ABIJ$ is a parallelogram with $AJ = BI = TU$.

The new parallelogram with $AC$ as one side and $RH$ the other, is equal, by our lemma.  $(ACDE) = (ACHR)$.  

Since $RAJ \parallel HCTU$ and $RA = HC = TU = AJ$, $(ATUJ) = (ACHR)$ for the same reason.  

So $(ATUJ) = (ACDE)$.

We use the same argument for $(BCFG) = (TBIU)$ on the right.

Then add the two results:  $(ACDE) + (BCFG) = (ABIJ)$.

$\square$.

Let $\angle ACB$ be a right angle, and let the parallelograms be squares.

\begin{center} \includegraphics [scale=0.30] {Pappus_pgram2.png} \end{center}

$\triangle DHC$ and $\triangle FHC$ are right triangles and they are congruent to $\triangle ABC$ by SAS.

So $\angle CHF = \angle CAT$ and we also have vertical angles, so $\triangle HFC \sim \triangle CTA$.

It follows that $\angle HFC = \angle CTA$ and both are right angles.

Further, $CH = TU = AB$, so $ABIJ$ is the square on $AB$.

It is easy to see that $(AEDC) = (ARHC) = (ATUJ)$, and the rest follows.

Thus Pappus' theorem becomes the Pythagorean theorem as a special case.

[ Reference:  George F. Simmons, \emph{Calculus Gems}. ]

\subsection*{proof without words}

%\begin{center} \includegraphics [scale=0.50] {pwow_pappus.png} \end{center}

\begin{center} \includegraphics [scale=0.40] {pappus_panels.png} \end{center}

[ Nelsen ]

\subsection*{question}

It might appear that this proof is not completely general because some choices for the parallelograms on sides $a$ and $b$ overlap at the top.  

Why is this not a problem?

\subsection*{Pythagoras yet again}

\label{sec:PProof_Leonardo}

Below (left) is a diagram for a proof of the Pythagorean theorem commonly ascribed to Leonardo da Vinci.  (Unfortunately, the attribution could not be verified and is most likely wrong.)

\begin{center} 
\includegraphics [scale=0.3] {Pyth_leo.png} 
\includegraphics [scale=0.3] {Pyth_leo2.png}
\includegraphics [scale=0.3] {Pyth_leo3.png} 
\end{center}

In the standard arrangement of I.47, connect the edges of the two smaller squares and form the polygon that cuts them both in half (red polygon).

Add a copy of the original right triangle to the bottom of the large square and join the vertices as shown to form the blue polygon.

Prove that the polygons shaded red and blue are congruent.

$\circ \ $  Show that $AN$ goes through the center of the large square and thus divides its area in half.  The other figures provide constructions that should help.

$\circ \ $  Show that $AN$ produces congruent parts of the right triangle.  (e.g. $\triangle AXG \cong \triangle NYL$).

$\circ \ $  Show that in addition to half the large square the blue polygon includes a whole copy of the right triangle.

$\circ \ $  Additionally, show that $AX$ bisects the right angle at $A$.  How does the central figure facilitate this?

\url{https://ar5iv.labs.arxiv.org/html/1311.0816}


\end{document}